\documentclass[11pt]{article}

% ============================================================
% GUÍA-RÚBRICA · EVALUACIÓN DE PARES DE EXPOSICIÓN
% ISR-401 Ingeniería de Requerimientos · SPA 2026-2027
% Aplica a: GA-SUM-02 (EXP-1, S4) y GA-SUM-06 (EXP-2, S13)
% Compilación: pdflatex -> pdflatex
% ============================================================

\usepackage[utf8]{inputenc}
\usepackage[T1]{fontenc}
\usepackage[scaled]{helvet}
\renewcommand{\familydefault}{\sfdefault}
\usepackage[a4paper,top=2.6cm,bottom=2.3cm,left=2cm,right=2cm,headheight=34pt]{geometry}
\usepackage{graphicx}
\usepackage{array}
\usepackage{tabularx}
\usepackage{longtable}
\usepackage{booktabs}
\usepackage{colortbl}
\usepackage{xcolor}
\usepackage{enumitem}
\usepackage{fancyhdr}
\usepackage{lastpage}
\usepackage[most]{tcolorbox}
\usepackage{amsmath}
\usepackage{caption}
\usepackage[hidelinks]{hyperref}
\usepackage{multirow}

% ---------- Textos en español (sin babel) ----------
\renewcommand{\tablename}{Tabla}
\renewcommand{\refname}{Referencias}

% ---------- Paleta institucional ----------
\definecolor{utegreen}{HTML}{0B7A3B}
\definecolor{utegreendk}{HTML}{075C2C}
\definecolor{midblue}{HTML}{1F5C8B}
\definecolor{gold}{HTML}{D4A017}
\definecolor{rowgray}{HTML}{EFEFEF}
\definecolor{boxgray}{HTML}{F4F6F4}
\definecolor{lightred}{HTML}{FBEAEA}
\definecolor{lightgold}{HTML}{FBF3DC}
\definecolor{lightblue}{HTML}{EAF1F7}
\definecolor{lightgreen}{HTML}{DCEFE3}

% ---------- Columnas personalizadas ----------
\newcolumntype{L}[1]{>{\raggedright\arraybackslash}p{#1}}
\newcolumntype{C}[1]{>{\centering\arraybackslash}p{#1}}
\newcolumntype{R}[1]{>{\raggedleft\arraybackslash}p{#1}}
\newcolumntype{Y}{>{\raggedright\arraybackslash}X}

% ---------- Encabezado y pie ----------
\pagestyle{fancy}
\fancyhf{}
\lhead{\footnotesize UTEQ $\cdot$ Ingeniería de Software $\cdot$ Ingeniería de Requerimientos (ISR-401)}
\rhead{\footnotesize Guía-Rúbrica de Evaluación de Pares · Exposición}
\lfoot{\footnotesize Prof. Gleiston C. Guerrero-Ulloa, PhD}
\rfoot{\footnotesize Página \thepage\ de \pageref{LastPage}}
\renewcommand{\headrulewidth}{0.6pt}
\renewcommand{\footrulewidth}{0.4pt}

\setlength{\parindent}{0pt}
\setlength{\parskip}{5pt}
\renewcommand{\arraystretch}{1.28}

% ---------- Atajos ----------
\newcommand{\thdr}[1]{\cellcolor{utegreen}\textcolor{white}{\textbf{#1}}}
\newcommand{\bhdr}[1]{\cellcolor{midblue}\textcolor{white}{\textbf{#1}}}

\newtcolorbox{cajaclave}[1][]{colback=boxgray,colframe=utegreen,boxrule=0.8pt,arc=2mm,fonttitle=\bfseries,#1}
\newtcolorbox{cajapiso}[1][]{colback=lightred,colframe=red!60!black,boxrule=1pt,arc=2mm,fonttitle=\bfseries,breakable,#1}
\newtcolorbox{cajaaviso}[1][]{colback=lightgold,colframe=gold!80!black,boxrule=0.8pt,arc=2mm,fonttitle=\bfseries,breakable,#1}

\begin{document}

% ============================================================
% PORTADA
% ============================================================
\thispagestyle{empty}
\noindent\includegraphics[width=\linewidth]{logoband.png}

\vspace{20pt}
\begin{center}
	{\large Facultad de Ciencias de la Computación}\\[2pt]
	{\large Carrera de Software}\\[30pt]
	{\LARGE\textbf{\textcolor{utegreendk}{GUÍA-RÚBRICA}}}\\[8pt]
	{\Large\textbf{Evaluación de Pares --- Informe de Exposición}}\\[6pt]
	{\large Ingeniería de Requerimientos (ISR-401)}\\[2pt]
	{\large Segundo Periodo Académico Ordinario (SPA) 2026--2027}
\end{center}

\vspace{20pt}
\begin{center}
	\begin{tabular}{L{0.30\linewidth}L{0.55\linewidth}}
		\textbf{Aplica a} & GA-SUM-02 (EXP-1, semana 4) y GA-SUM-06 (EXP-2, semana 13)\\
		\textbf{Tipo de actividad} & GA-EXP (Gestión de Aula --- Exposición)\\
		\textbf{Rol del estudiante} & Evaluador individual de los equipos que NO son el suyo\\
		\textbf{Docente} & Gleiston Cicerón Guerrero Ulloa, PhD\\
		\textbf{Correo institucional} & \texttt{gguerrero@uteq.edu.ec}\\
	\end{tabular}
\end{center}

\vspace{18pt}
\noindent\textcolor{gold}{\rule{\linewidth}{1.5pt}}
\begin{center}
	\small Este documento describe qué evalúa el par observador, cómo lo registra y con qué criterios se califica el informe resultante. La nota de la \emph{exposición} es la que asigna el docente; la nota de este \emph{informe de evaluación de pares} se registra de forma independiente como parte de la participación en clase (GA).
\end{center}

\newpage

% ============================================================
\section{\textcolor{utegreendk}{Propósito y contexto}}
% ============================================================

Cada vez que un equipo expone, los integrantes de los demás equipos actúan como evaluadores pares: observan, registran evidencias y emiten un juicio fundamentado sobre la calidad de la exposición. Este ejercicio tiene dos propósitos: (1) desarrollar en el observador la capacidad de evaluar críticamente el trabajo ajeno con criterios técnicos objetivos y medibles, y (2) proporcionar al equipo expositor retroalimentación adicional a la del docente.

El \textbf{informe de evaluación de pares} es el documento escrito que el estudiante-evaluador entrega al final de la sesión de exposiciones. Cada estudiante entrega \textbf{un informe por cada equipo observado} (no uno propio). No se evalúa al equipo al que pertenece el evaluador.

\begin{cajaclave}[title=Regla de exclusión]
	El estudiante \textbf{no evalúa} a los integrantes de su propio equipo de exposición. Si el grupo tiene cuatro equipos (A, B, C, D) y el evaluador es del equipo A, entrega informes para los equipos B, C y D.
\end{cajaclave}

\begin{cajaclave}[title=Qué NO debe hacer el evaluador]
	\begin{itemize}[leftmargin=14pt,itemsep=2pt]
		\item No filma ni graba audio o video a los expositores.
		\item No interrumpe la exposición para hacer preguntas (salvo en el turno de preguntas habilitado por el docente).
		\item No consulta con sus compañeros durante la observación para ponerse de acuerdo en las puntuaciones.
		\item No copia el informe de otro evaluador.
	\end{itemize}
\end{cajaclave}

% ============================================================
\section{\textcolor{utegreendk}{Estructura del informe de evaluación de pares}}
% ============================================================

El informe se entrega en papel o en formato digital (PDF) al docente al finalizar la sesión, según se indique en cada actividad. Su extensión es de \textbf{una a dos páginas por equipo evaluado}. Debe contener las secciones descritas a continuación.

\subsection{\textcolor{midblue}{Sección 1 --- Identificación}}

\noindent\begin{tabularx}{\linewidth}{|L{0.38\linewidth}|Y|}
	\hline
	\bhdr{Campo} & \bhdr{Descripción}\\
	\hline
	Nombre completo del evaluador & Quién escribe el informe.\\
	\hline
	\rowcolor{rowgray} Equipo del evaluador & Nombre o número del equipo al que pertenece.\\
	\hline
	Equipo evaluado & Nombre o número del equipo que expuso.\\
	\hline
	\rowcolor{rowgray} Tema de la exposición & Título exacto que el equipo expositor declaró al inicio.\\
	\hline
	Actividad & GA-SUM-02 (EXP-1) o GA-SUM-06 (EXP-2).\\
	\hline
	\rowcolor{rowgray} Fecha & Día en que se realizó la exposición.\\
	\hline
\end{tabularx}

\subsection{\textcolor{midblue}{Sección 2 --- Registro de indicadores observados}}

Para cada criterio de la Tabla~\ref{tab:criterios}, el evaluador registra:
\begin{enumerate}[leftmargin=18pt,itemsep=3pt]
	\item \textbf{Nivel asignado} (0, 1, 2, 3 o 4) según los descriptores de la Tabla~\ref{tab:niveles}.
	\item \textbf{Evidencia concreta}: al menos una oración que describa qué observó durante la exposición y que justifique el nivel asignado. La evidencia debe ser observable y medible (no opiniones como «fue bueno»); debe citar, por ejemplo, el número de diapositiva, un término técnico usado o no usado, el tiempo observado, una pregunta que no pudo responder, etc.
\end{enumerate}

\subsection{\textcolor{midblue}{Sección 3 --- Nota calculada}}

El evaluador calcula la nota con la fórmula de la sección~\ref{sec:formula} y la registra en la tabla de criterios.

\subsection{\textcolor{midblue}{Sección 4 --- Retroalimentación constructiva}}

Dos o tres frases dirigidas al equipo expositor: qué fue lo más valioso de la presentación y qué aspecto concreto mejoraría más su calidad técnica. Sin adjetivos vacíos («fue interesante»): cada frase debe identificar un elemento específico y medible.

% ============================================================
\section{\textcolor{utegreendk}{Criterios de evaluación de la exposición}}
% ============================================================

Los criterios de la Tabla~\ref{tab:criterios} describen lo que el evaluador par observa en la exposición del equipo. Cada criterio tiene un peso que refleja su importancia en la calidad de una exposición técnica de IR. Los pesos siguen la distribución GA-EXP del encuadre: D1 = 10, D2 = 20, D3 = 15, D4 = 20, D5 = 25, D6 = 10.

\begin{cajaaviso}[title=Indicadores observables y medibles]
	Cada descriptor incluye indicadores concretos: tiempo, número de requisitos, presencia o ausencia de un elemento, coherencia entre lo dicho y lo mostrado en pantalla. Asigne el nivel que mejor refleje lo observado; si un indicador intermedio no se da en ningún nivel exactamente, aplique el más cercano y justifíquelo en la evidencia.
\end{cajaaviso}

{\footnotesize\setlength{\tabcolsep}{3pt}
\begin{longtable}{|C{0.45cm}|L{2.4cm}|C{0.65cm}|L{3.1cm}|L{3.1cm}|L{3.1cm}|L{2.4cm}|}
	\caption{Criterios de evaluación de la exposición para el evaluador par.}\label{tab:criterios}\\
	\hline
	\thdr{\#} & \thdr{Criterio} & \thdr{Peso} & \thdr{4 --- Excelente} & \thdr{3 --- Bueno} & \thdr{2/1 --- Aceptable/Insuficiente} & \thdr{0 --- Ausente}\\
	\hline
	\endfirsthead
	\hline
	\thdr{\#} & \thdr{Criterio} & \thdr{Peso} & \thdr{4 --- Excelente} & \thdr{3 --- Bueno} & \thdr{2/1 --- Aceptable/Insuficiente} & \thdr{0 --- Ausente}\\
	\hline
	\endhead
	\hline
	\multicolumn{7}{r}{\footnotesize\textit{Continúa en la página siguiente\ldots}}\\
	\endfoot
	\hline
	\endlastfoot
	D1 & \textbf{Cumplimiento de la guía} (estructura y contenido mínimo) & 10 & Trata todos los puntos pedidos en la guía de la actividad: técnicas/herramientas $\geq$ las indicadas, comparación y recomendación presentes, demostración realizada. & Falta un punto menor (p.\,ej., la recomendación es superficial) pero los demás están presentes. & \textbf{(2)} Falta más de un punto o uno de los requeridos (comparación o demostración) está ausente. \textbf{(1)} Solo expone un tema o la mitad de los puntos. & No expone o el tema no corresponde a la guía.\\
	\hline
	\rowcolor{rowgray} D2 & \textbf{Fundamentación conceptual} (precisión técnica y uso de normas) & 20 & Usa términos técnicos de IR con exactitud ($\geq$3 conceptos del encuadre nombrados correctamente); las afirmaciones sobre las técnicas/herramientas son verídicas y coherentes con IEEE 29148, SWEBOK o IREB. & Conceptos mayoritariamente correctos; $\leq$1 imprecisión que no invalida el contenido. & \textbf{(2)} Dos o más imprecisiones o términos usados de forma incorrecta. \textbf{(1)} Errores conceptuales graves o confusión sistemática entre conceptos. & No hay fundamentación o el tema no pertenece a IR.\\
	\hline
	D3 & \textbf{Desarrollo técnico} (profundidad del contenido técnico presentado) & 15 & Compara $\geq$2 técnicas o herramientas con $\geq$3 criterios explícitos y medibles (funcionalidad, madurez, integración, costo, etc.); los datos son verificables (fuente citada en diapositiva). & Comparación con $\geq$2 criterios; algunos datos sin fuente, pero el análisis es correcto. & \textbf{(2)} Comparación descriptiva sin criterios explícitos. \textbf{(1)} Solo enumera características sin comparar ni analizar. & No hay contenido técnico evaluable.\\
	\hline
	\rowcolor{rowgray} D4 & \textbf{Análisis crítico y recomendación} (juicio fundamentado) & 20 & La recomendación final (técnica o herramienta) está argumentada con $\geq$2 razones específicas del contexto del PFC u otro caso concreto; reconoce limitaciones. & Recomendación presente con $\geq$1 argumento; limitaciones no mencionadas o superficiales. & \textbf{(2)} Recomendación sin argumento o genérica. \textbf{(1)} No hay recomendación o es contradictoria con el análisis. & Ausente.\\
	\hline
	D5 & \textbf{Comunicación oral y visual} (claridad, tiempo, respuestas) & 25 & Exposición fluida sin leer diapositivas; tiempo dentro del rango $\pm$2 min; responde $\geq$1 pregunta técnica con precisión; diapositivas legibles (fuente $\geq$18 pt visible desde el fondo del aula). & Uno de los cuatro indicadores cumplido parcialmente (pequeña lectura, tiempo $\pm$3 min, respuesta con imprecisión menor, una diapositiva ilegible). & \textbf{(2)} Dos indicadores parciales. \textbf{(1)} Tres o más indicadores parciales o no responde ninguna pregunta. & No expone o la exposición es ininteligible.\\
	\hline
	\rowcolor{rowgray} D6 & \textbf{Fuentes y créditos} (atribución del contenido) & 10 & Cada dato, captura o estadística tiene fuente visible en la diapositiva correspondiente; las fuentes son verificables (autor, año, URL o norma). & Mayoría de fuentes presentes; $\leq$2 datos sin fuente pero el contenido es correcto. & \textbf{(2)} Varias fuentes ausentes. \textbf{(1)} Solo una o ninguna fuente visible en todo el juego de diapositivas. & No hay ninguna fuente o el contenido es copiado sin atribución.\\
	\hline
\end{longtable}
}

% ============================================================
\section{\textcolor{utegreendk}{Niveles de desempeño}}\label{sec:niveles}
% ============================================================

\noindent{\small\begin{tabularx}{\linewidth}{|C{0.07\linewidth}|L{0.18\linewidth}|Y|}
	\hline
	\thdr{Nivel} & \thdr{Denominación} & \thdr{Descripción operacional para el evaluador par}\\
	\hline
	4 & Excelente & Cumple íntegramente todos los indicadores del descriptor; no se detecta ninguna omisión ni imprecisión en lo observado.\\
	\hline
	\rowcolor{rowgray} 3 & Bueno & Cumple la mayor parte; hay una deficiencia menor que no afecta la comprensión ni la validez técnica.\\
	\hline
	2 & Aceptable & Cumple la mitad aproximadamente; hay deficiencias que limitan pero no anulan la comprensión.\\
	\hline
	\rowcolor{rowgray} 1 & Insuficiente & Cumple menos de la mitad o hay errores graves que comprometen la comprensión técnica.\\
	\hline
	0 & Ausente & El criterio no se puede evaluar porque el elemento no se presentó o no fue observable.\\
	\hline
\end{tabularx}}\label{tab:niveles}

% ============================================================
\section{\textcolor{utegreendk}{Cálculo de la nota sugerida}}\label{sec:formula}
% ============================================================

El evaluador par calcula una \textbf{nota sugerida} que el docente toma como insumo (no como nota definitiva). La nota definitiva de la exposición la asigna el docente.

\[
	\text{Nota sugerida} = \frac{\displaystyle\sum_{i=1}^{6} w_i \times n_i}{4} \times \frac{10}{100}
\]

donde $w_i$ es el peso de la dimensión $i$ y $n_i$ el nivel asignado (0 a 4). Los pesos son: D1 = 10, D2 = 20, D3 = 15, D4 = 20, D5 = 25, D6 = 10.

\begin{cajaclave}[title=Ejemplo de cálculo]
	Niveles observados: D1 = 4, D2 = 3, D3 = 3, D4 = 3, D5 = 4, D6 = 2.\\[4pt]
	$\text{Nota sugerida} = \dfrac{10 \cdot 4 + 20 \cdot 3 + 15 \cdot 3 + 20 \cdot 3 + 25 \cdot 4 + 10 \cdot 2}{4} \times \dfrac{10}{100} = \dfrac{305}{4} \times 0{,}1 = 7{,}6$
\end{cajaclave}

% ============================================================
\section{\textcolor{utegreendk}{Criterios de calificación del informe de evaluación de pares}}
% ============================================================

El \emph{informe} que entrega el evaluador es en sí mismo una actividad evaluada. Se califica sobre 10 con los criterios de la Tabla~\ref{tab:rubrica_informe}.

\begin{cajapiso}[title={Criterios de piso del informe: incumplimiento implica CERO en el informe}]
	\begin{enumerate}[label=\textbf{P\arabic*.},leftmargin=24pt,itemsep=3pt]
		\item \textbf{Entrega al docente al finalizar la sesión.} No se acepta entrega posterior.
		\item \textbf{Un informe por cada equipo observado} (excluido el propio). Si faltan informes de alguno de los equipos observados, la nota del informe es CERO.
		\item \textbf{Evidencia concreta por criterio.} Si más de dos criterios carecen de evidencia observable (solo nivel sin texto), la nota es CERO.
		\item \textbf{Sin copia de otro evaluador.} Informes con evidencias idénticas o muy similares (misma redacción, mismos ejemplos) reciben CERO en ambos evaluadores.
	\end{enumerate}
\end{cajapiso}

{\footnotesize\setlength{\tabcolsep}{3pt}
\begin{longtable}{|L{2.6cm}|C{0.85cm}|L{3.3cm}|L{3.0cm}|L{3.3cm}|L{2.2cm}|}
	\caption{Rúbrica de calificación del informe de evaluación de pares.}\label{tab:rubrica_informe}\\
	\hline
	\thdr{Criterio del informe} & \thdr{Peso} & \thdr{4 --- Excelente} & \thdr{3 --- Bueno} & \thdr{2/1 --- Acept./Insuf.} & \thdr{0 --- Ausente}\\
	\hline
	\endfirsthead
	\hline
	\thdr{Criterio del informe} & \thdr{Peso} & \thdr{4 --- Excelente} & \thdr{3 --- Bueno} & \thdr{2/1 --- Acept./Insuf.} & \thdr{0 --- Ausente}\\
	\hline
	\endhead
	\hline
	\multicolumn{6}{r}{\footnotesize\textit{Continúa\ldots}}\\
	\endfoot
	\hline
	\endlastfoot
	\textbf{I1. Completitud} (todos los campos y criterios completos) & 20 & Identificación completa; los 6 criterios con nivel y evidencia; nota calculada; retroalimentación presente. & Falta un campo menor de identificación o un criterio sin evidencia. & \textbf{(2)} Dos criterios sin evidencia o falta la sección de retroalimentación. \textbf{(1)} Tres o más criterios sin evidencia. & No entregó o entregó una hoja en blanco.\\
	\hline
	\rowcolor{rowgray} \textbf{I2. Objetividad y base observable} (evidencias concretas, no opiniones) & 30 & Cada nivel está respaldado por al menos una observación concreta y medible (diapositiva, término, tiempo, pregunta); sin afirmaciones como «fue bueno». & Al menos 4 de 6 criterios tienen evidencia concreta; los otros son superficiales. & \textbf{(2)} Solo 2--3 criterios con evidencia concreta. \textbf{(1)} Evidencias todas genéricas u opinativas. & Sin ninguna evidencia observable.\\
	\hline
	\textbf{I3. Coherencia entre nivel y evidencia} (el nivel asignado se justifica con lo descrito) & 25 & En todos los criterios el nivel asignado es coherente con la evidencia descrita (no asigna 4 con evidencia de 2, ni viceversa). & En $\leq$1 criterio el nivel y la evidencia no son del todo coherentes. & \textbf{(2)} En 2--3 criterios hay incoherencia. \textbf{(1)} En 4 o más criterios el nivel no concuerda con lo descrito. & Sin coherencia verificable.\\
	\hline
	\rowcolor{rowgray} \textbf{I4. Cálculo correcto} (nota sugerida calculada correctamente) & 10 & La nota sugerida se calcula con la fórmula correcta, los pesos correctos y sin error aritmético. & Error aritmético menor ($\leq$0,2 puntos de diferencia) pero fórmula y pesos correctos. & \textbf{(2)} Error en un peso o en la fórmula. \textbf{(1)} Nota sugerida ausente o calculada sin fórmula. & Sin nota calculada.\\
	\hline
	\textbf{I5. Retroalimentación constructiva} (específica y accionable) & 15 & Identifica un punto fuerte y uno de mejora concretos, con referencia a un elemento específico de la exposición (no genérico). & Uno de los dos (fortaleza o mejora) es concreto; el otro es genérico. & \textbf{(2)} Ambos son genéricos pero presentes. \textbf{(1)} Solo uno de los dos presente. & Ausente.\\
	\hline
\end{longtable}
}

Fórmula del informe:
\[
	\text{Nota del informe} = \frac{20 \cdot n_{I1} + 30 \cdot n_{I2} + 25 \cdot n_{I3} + 10 \cdot n_{I4} + 15 \cdot n_{I5}}{4} \times \frac{10}{100}
\]

% ============================================================
\section{\textcolor{utegreendk}{Plantilla de entrega (llenar a mano o en texto)}}
% ============================================================

\begin{tcolorbox}[colback=lightblue,colframe=midblue,boxrule=0.8pt,arc=2mm,title=Informe de Evaluación de Pares --- Exposición ISR-401]
	\small
	\noindent\textbf{Evaluador:}~\underline{\hspace{5.5cm}}\hfill\textbf{Equipo del evaluador:}~\underline{\hspace{3cm}}\\[5pt]
	\noindent\textbf{Equipo evaluado:}~\underline{\hspace{4.5cm}}\hfill\textbf{Actividad:}~\underline{\hspace{2.2cm}}\hfill\textbf{Fecha:}~\underline{\hspace{2cm}}\\[5pt]
	\noindent\textbf{Tema de la exposición:}~\underline{\hspace{10.2cm}}\\[8pt]

	\renewcommand{\arraystretch}{1.5}
	\begin{tabularx}{\linewidth}{|C{0.7cm}|L{2.8cm}|C{0.9cm}|Y|}
		\hline
		\textbf{D} & \textbf{Criterio} & \textbf{Nivel} & \textbf{Evidencia concreta observada}\\
		\hline
		D1 & Cumplimiento guía & & \\[12pt]
		\hline
		D2 & Fundamentación conceptual & & \\[12pt]
		\hline
		D3 & Desarrollo técnico & & \\[12pt]
		\hline
		D4 & Análisis y recomendación & & \\[12pt]
		\hline
		D5 & Comunicación oral y visual & & \\[12pt]
		\hline
		D6 & Fuentes y créditos & & \\[12pt]
		\hline
		\multicolumn{2}{|r|}{\textbf{Nota sugerida calculada:}} & \multicolumn{2}{l|}{\hspace{0.5cm}/10}\\
		\hline
	\end{tabularx}

	\vspace{6pt}
	\textbf{Retroalimentación al equipo expositor:}\\[4pt]
	Punto fuerte: \underline{\hspace{12cm}}\\[8pt]
	Aspecto a mejorar: \underline{\hspace{11.5cm}}
\end{tcolorbox}

\end{document}
